\begin{Example}
\label{ex:stokes-dirichlet-complementing-condition}
Consider the Stokes system $\Delta\bm{u}-\nabla p=0$, $\operatorname{div}\bm{u}=0$, where $\bm{u}\in\mathbb{R}^3$, $p\in\mathbb{R}$, in the half-space $z>0$, with the Dirichlet boundary condition $\bm{u}=0$ on $z=0$. As we showed in Chapter~\ref{chap:characteristics}, the system is elliptic, and of order $6$. Assume now that we have a solution $\bm{u}(x,y,z)=\exp(i\xi_1x+i\xi_2y)\bm{v}(z)$, $p(x,y,z)=\exp(i\xi_1x+i\xi_2y)q(z)$, where $\bm{v}$ and $q$ tend to zero exponentially as $z\to\infty$. Let $\Sigma=(0,L_1)\times(0,L_2)$, where $L_i$ is a multiple of $2\pi/\xi_i$ if $\xi_i\ne0$ and arbitrary if $\xi_i=0$. We find
\begin{equation}\label{eq:stokes-complementing-periodic-energy-identity}
0=\int_{\Sigma\times\mathbb{R}^+}(\Delta\bm{u}-\nabla p)\cdot\bm{u}\,\mathrm{d}x\,\mathrm{d}y\,\mathrm{d}z=-\int_{\Sigma\times\mathbb{R}^+}|\nabla\bm{u}|^2\,\mathrm{d}x\,\mathrm{d}y\,\mathrm{d}z.
\end{equation}
This implies that $\bm{u}$ is constant, which is compatible with our assumptions only if $\bm{u}=0$. It easily follows that $p$ is also zero. Hence the Stokes system with Dirichlet boundary conditions satisfies the complementing condition.
\end{Example}

\begin{Example}
\label{ex:biharmonic-noncomplementing-boundary-condition}
Consider the biharmonic equation $\Delta\Delta u=0$ in the half-plane $y>0$ with boundary conditions $\Delta u=\frac{\partial}{\partial y}\Delta u=0$ on the line $y=0$. For every $\xi\in\mathbb{R}$, the function $u(x,y)=\exp(i\xi x-|\xi|y)$ is a solution. Hence the complementing condition does not hold.
\end{Example}

The complementing condition or its failure is not always as easy to verify as in the preceding examples. However, it can always be reduced to a purely algebraic problem. If we insert the ansatz \eqref{eq:complementing-decaying-tangential-mode} into \eqref{eq:complementing-frozen-principal-system}, we obtain a system of ODEs for $\bm{v}(\eta)$, which can as usual be solved by the ansatz $\bm{v}(\eta)=\exp(\lambda\eta)\bm{v}_0$. Ellipticity means that no roots $\lambda$ are imaginary, and if the coefficients of our system are real, then an equal number of roots must have positive and negative real parts. Let $\lambda_l^+(\bm{x}_0,\bm{\xi})$ denote the roots with positive real part. Then one obtains $m$ linearly independent solutions of \eqref{eq:complementing-frozen-principal-system} in the form $\exp(i\bm{\xi}\cdot(\bm{x}-\bm{x}_0)+\lambda_l^+\eta)\bm{u}_l$ (in the usual way, this may need to be modified by including powers of $\eta$ if there are repeated roots). It remains to be checked if any linear combination of these solutions satisfies the boundary conditions, which is a purely algebraic problem. For equivalent algebraic characterizations of the complementing conditions we refer to the literature, see \cite{ADN2}.
